\chapter{Analyse vectorielle dans $\Rn$}
\lecture{7 septembre}
\begin{definition}[Définition de base]
	Soit $U \subset \Rn$ un \textbf{domaine} \label{def: Domaine}, c'est-à-dire un sous-ensemble \emph{ouvert} et \emph{connexe}, on appelle
	\begin{enumerate}
		\item $f \in C^k(U, \R), ~k \geqslant 0$ un \textbf{champ scalaire}. \label{def: Champ scalaire}
		\item $\bm f \in C^k(U, \Rn)$ un \textbf{champ vectoriel}. \label{def: Champ vectoriel}
	\end{enumerate}
\end{definition}

\begin{definition}[Gradient] \label{def: Gradient}
	Soit $g \in C^1(U, \R)$ un champ scalaire. Alors on apelle
	\[
	\bm f = \nabla g
	\]
	le \textbf{gradient} de $g$. Dans ce cas, on dit que $\bm f$ dérive d'un champ scalaire $\bm f$.
\end{definition}

 \section{Notions d'intégration}
 
 \begin{definition}[chemin] \label{def: Chemin}
 	Un \textbf{chemin} est une application continue 
 	\[
 	\gamma : I \to U
 	\]
 	où $I = [s,t], \quad s < t \in \R$, et tel qu'il existe une partition de $I = \bigcup_{j=1}^N I_j$ où les $I_j$ sont des intervalles fermés et d'intérieurs disjoints, tel que les restrictions
 	\[
 	\gamma_{\mid I_j} \in C^1(I_j, U)
 	\]
 	autrement dit $\gamma$ est \emph{continue par morceaux}. De plus on dira que $\gamma$ est un \textbf{lacet} ou \textbf{chemin fermé} si $\gamma(s) = \gamma(t)$.
 \end{definition}
 
 \subsection{Intégrale curviligne}

 \begin{definition}[Intégrale curviligne] \label{def: Intégrale curviligne}
 	Soit $f \in C^0(U, \R)$ un champ scalaire, et soit $\gamma : I \to U$ un chemin, alors on définit l'\textbf{intégrale curviligne} de $f$ le long de $\gamma$ par
 	\[
 	\int_\gamma f \, \mathrm ds = \int_I f(\gamma(t)) \, \norme{\gamma'(t)} \, \mathrm dt =: \sum_{j=1}^N \int_{I_j} f(\gamma(t)) \, \norme{\gamma'(t)} \, \mathrm dt
 	\]
 	où $\norme{ \gamma'(t)} = \sqrt{\sum_{i=1}^n \gamma_i'(t)^2}$.
 \end{definition}
 
 \begin{definition}[Reparamétrisation d'un chemin] \label{def: Reparamétrisation d'un chemin}
 	Soit $\gamma : [s,t] \to U$ un chemin et soit
 	\[
 	\varphi : [\tilde s, \tilde t] \to [s,t]
 	\]
 	un difféomorphisme de classe $C^1$ tel que $\varphi' > 0$. Alors la composition $\gamma \circ \varphi : [\tilde s, \tilde t] \to U$ est encore un chemin.
 \end{definition}
 
 \begin{theorem}[Invariance par reparamétrisation] \label{thm: Invariance par reparamétrisation}
 	Soit $U \subset \Rn$ un domaine, $f \in C^0(U)$ un champ scalaire, $\gamma : I \to U$ et $\gamma \circ \varphi : \tilde I \to U$ une reparamétrisation. Alors
 	\[
 	\int_{\gamma \circ \varphi} f \, \mathrm ds = \int_\gamma f \, \mathrm ds
 	\]
\end{theorem}
\prf{Soit $\tilde I = [\tilde s, \tilde t]$, $I =[s,t]$. Alors
	\[
	\int_{\gamma \circ \varphi} f \, \mathrm ds = \int_{\tilde I} f(\gamma \circ \varphi \, (x)) \, \norme{(\gamma \circ \varphi)'(x)} \, \mathrm dx = \int_{\tilde s}^{\tilde t} f(\gamma \circ \varphi(x)) \, \norme{\gamma'\circ \varphi(x) \, \varphi'(x)} \, \mathrm dx = \int_{\tilde s}^{\tilde t} f(\gamma \circ \varphi(x)) \varphi'(x) \, \norme{\gamma' \circ \varphi(x)} \, \mathrm dx
	\]
	et donc par changement de variable on a bien, 
	\[
	\int_{\gamma \circ \varphi} f \, \mathrm ds = \int_{\varphi(\tilde s)}^{\varphi(\tilde t)} f(\gamma(y)) \, \norme{\gamma'(y)} \, \mathrm dy	 = \int_\gamma f \, \mathrm ds	
	\]
	d'où le résultat.
}

\subsection{Circulation d'un champ vectoriel}

\begin{definition}[Circulation] \label{def: Circulation}
	 Soit $U \subset \Rn$ un domaine et $\bm f \in C^0(U, \Rn)$ un champ vectoriel et $\gamma : I \to U$ un chemin, alors on définit la \textbf{circulation} du champ vectoriel le long de $\gamma$ par
	\[
	\int_{\gamma} \bm f \cdot \mathrm d \bm s = \int_I \bm f(\gamma(t)) \cdot \bm \gamma'(t) \, \mathrm dt
	\]
	où $\cdot$ désigne le produit scalaire euclidien.
\end{definition}

\begin{remarque} L'intégrale est invariante par reparamétrisation. Mais si on change le sens de parcours, alors le signe change. Par exemple, soit $\gamma : [0,1] \to U$ un chemin et $\varphi : [0,1] \to [0,1], ~\varphi(t) = 1-t$ alors on trouve bien
	\[
	 \int_{\gamma \circ \varphi} \bm f \cdot \, \mathrm \bm s = - \int_{\gamma} \bm f \cdot \, \mathrm d \bm s
	\]
	alors on peut noter plus simplement $\gamma \circ \varphi = - \gamma$.
\end{remarque}

\begin{definition}[Concaténation de chemin] \label{def: Concaténation de chemin}
	Soit $\gamma : [a,b] \to U$ et $\tilde \gamma : [b, c] \to U$ des chemins tels que $\gamma(b) = \tilde \gamma(b)$, alors on définit la \textbf{concaténation} $\gamma\oplus \tilde \gamma$ par 
	\[
	\begin{array}{cccc}
		\gamma \oplus \tilde \gamma : & [a,c] & \longrightarrow & U \\
										& t & \longmapsto & 
										\begin{cases}
											\gamma(u), & u \in [a,b] \\
											\tilde\gamma(u), & u \in [b,c]
										\end{cases}
	\end{array}
	\]
	Il est facile de voir que $\gamma \oplus \tilde \gamma$ est bien $C^1$ par morceaux et donc il s'agit bien d'un chemin.
\end{definition}

On se demande s'il existe des champs vectoriels $\bm f$ tels que $\int_\gamma \bm f \cdot \mathrm d \bm s$ ne dépend que du point initial et du point final de $\gamma$.

\begin{lemme}
	Soit $U \subset \Rn$ un domaine, $\bm f \in C^k(U)$ un champ vectoriel. Alors $\bm f$ à la propriétés que $\int_\gamma \bm f \cdot \, \mathrm d \bm s$ ne dépend que des extrémités des chemins si et seulement si $\int_\gamma \bm f \cdot \, \mathrm d \bm s  = 0$ pour chaque lacet $\gamma : I \to U$.
\end{lemme}

\prf{Si $\int_\gamma \bm f \cdot \, \mathrm d \bm s$ ne dépend que des extrémités, on considère un lacet $\gamma$ arbitraire. Alors on sait qu'il existe $\gamma_1$ et $\gamma_2$ tel que $\gamma = \gamma_1 \oplus -\gamma_2$  et donc a 
	\[
	\int_\gamma \bm f \cdot \mathrm d \bm s = \int_{\gamma_1 \oplus -\gamma_2} \bm f \cdot \mathrm d \bm s = \int_{\gamma_1} \bm f \cdot \mathrm d \bm s - \int_{\gamma_2} \bm f \cdot \mathrm d \bm s = 0
	\]	
	Et inversement, alors si $\int_\gamma \bm f \cdot \, \mathrm d \bm s = 0$ pour tout lacet $\gamma : I  \to U$m  et si $\gamma_1, \gamma_2$ lient $\bm p$ et $\bm q$, alors on peut introduire le lacet qui la concaténation de $\gamma_1$ et $-\gamma_2$ et donc
	\[
	\int_{\gamma_1} \bm f \cdot \, \mathrm d \bm s - \int_{\gamma_2} \bm f \cdot \, \mathrm d \bm s = \int_{\gamma_1 \oplus - \gamma_2} \bm f \cdot \, \mathrm d \bm s = 0
	\]
	et donc on a bien la propriété voulue.
}


\lecture{10 septembre}
7
\begin{theorem}
	Soit $U \subset \Rn$ un domaine, et $\bm f \in C^k(U)$ un champ vectoriel, alors $\int_\gamma \bm f\cdot \, \mathrm d \bm s = 0$ pour tout lacet $\gamma$ si et seulement si $\bm f$ dérive d'un potentiel, donc s'il existe $g \in C^{k+1}(U, \R)$ tel que $\bm f = \nabla g$.
\end{theorem}



\prf{Supposons que $\bm f = \nabla g$, et que $\gamma : [0,1] \to U$ est un lacet, alors
\[
\int_\gamma \bm f \cdot \mathrm d \bm s = \int_0^1 \nabla g(\gamma(t)) \cdot \gamma'(t) \, \mathrm dt = \int_0^1 \frac{\mathrm d}{\mathrm dt} g(\gamma(t)) \, \mathrm dt = g(\gamma(1)) - g(\gamma(0))
\]
mais comme $\gamma$ est un lacet alors $\gamma(1) = \gamma(0)$ et donc on a bien
\[
\int_\gamma \bm f \cdot \mathrm d \bm s = 0
\]

Réciproquement supposons que $\int_\gamma \bm f \cdot \mathrm d \bm s = 0$ pour tout lacet, alors il faut exhiber $g \in C^1(U)$ avec $\bm f = \nabla g$. Par connexité de $U$, on considère $p \in U$, et on pose
\[
g(x) = \int_\gamma \bm f \cdot \mathrm d \bm s
\]
où $\gamma$ est un chemin arbitraire en $U$ qui relie $p$ à $x$. D'après  les hypothèse, alors $g$ est bien défini. Il est facile de voir que $g$ est $C^0$, il reste donc à montrer que $g$ est $C^1$ et qu'on a bien $\nabla g = \bm f$. Or pour montrer que $\nabla g = f$, et donc à fortiori que $g \in C^1(U)$ on calcule les dérivées partielles de $g$. Alors on a 
\[
\frac{\partial g}{\partial x_j} (x) = \lim_{h \to 0} \frac{g(x + h e_j) - g(x)}{h} 
\]
or 
\[
g(x + h e_j) = \int_{\tilde \gamma} \bm f \cdot \mathrm d\bm s
\]
avec $\tilde \gamma $ est un chemin arbitraire qui lie $p$ à $x + he_j$. Or par concaténation on peut écrire, $\tilde \gamma = \gamma \oplus \gamma_1$ où $\gamma_1 : [0,1] \to U, \quad \gamma_1(t) = x  + t h e_j$
Et donc par hypothèse
\[
g(x +he_j) = \int_{\gamma \oplus \gamma_1} \bm f \cdot \mathrm d \bm s = \int_\gamma \bm f \cdot \mathrm d \bm s + \int_{\gamma_1} \bm f \cdot \mathrm d \bm s
\]
Et donc on a
\[
g(x + h e_j) - g(x) = \int_{\gamma_1} \bm f \cdot \mathrm d \bm s = \int_0^1 \bm f(x + th e_j) \cdot h e_j \, \mathrm dt = h \int_0^1 f_j(x + the_j) \, \mathrm dt = h \int_0^1 f_j(x) \, \mathrm dt + h \int_0^1 r_j(t,h,x) \, \mathrm dt
\]
où $r_j(t,h,x) = f_j(x + the_j) - f_j(x) \longrightarrow 0$ et donc $h \int_0^1 r_j(t,h,x) \, \mathrm dt = o(h)$ et donc 
\[
\frac{\partial g}{\partial x_j} (x) = \lim_{h \to 0} \frac{g(x + h e_j) - g(x)}{h} = f_j(x)
\]
Ainsi les dérivées partielles existent, sont continues, et $g$ est continue, ainsi $g \in C^1(U)$ et $\nabla g = \bm f$.
}

\begin{remarque}
	On voit que si $\bm f = \nabla g$ avec $g \in C^2(U)$ alors
	\[
	\frac{\partial f_j}{\partial k} = \frac{\partial^2 g}{\partial x_k \partial x_j} = \frac{\partial^2 g}{\partial x_j \partial x_k} = \frac{\partial f_k}{\partial x_j}
	\]
	Ainsi on condition nécessaire pour que $\bm f \in C^1(U, \Rn)$ dérive d'un potentiel est alors que
	\[
	\forall i,j \in \entiers, \quad \frac{\partial f_j}{\partial x_i} - \frac{\partial f_i}{\partial x_j}
	\]
	Cependant dans le cadre de ce cours, on va se restreindre au cas $n=2$ ou $n=3$, alors on a la condition nécessaire
	\[
	\bm f = \nabla g \implies \rot \bm f = 0
	\]
\end{remarque}

\subsection{Concepts topologiques}

\begin{definition}[Homotopie] \label{def: Homotopie}
	Soit $U \subset \Rn$ un domaine, et soient $\gamma_1, \gamma_2 : [0,1] \to U$ deux lacets dans $U$ possédant les mêmes points d'extrémité. Alors on dit que $\gamma_1$ est \textbf{homotope} à $\gamma_2$ s'il existe une application continue
	\[
	\Gamma: [0,1] \times [0,1] \to U
	\]
	telle que $\Gamma(0,t) =  \gamma_1(t)$ et $\Gamma(1,t) = \gamma_2(t)$ et pour chaque $s \in [0,1]$ fixé, $\Gamma(s,t)$ est un chemin qui lie les mêmes point d'extrémité que $\gamma_1$ et $\gamma_2$. La définition est analogue pour les lacets.
\end{definition}


\begin{definition}[Contractibilité et simplement connexe] \label{def: Contractibilité et simplement connexe}
	Soit $U \subset \Rn$ un domaine, et $\gamma : [0,1] \to U$ un lacet, alors on dit que $\gamma$ est \textbf{contractible} pourvu qu'il soit homotope à un lacet constant en $U$. En plus si $U \subset \Rn$ dont tous les lacets sont contractibles, alors on dira que $U$ est \textbf{simplement connexe}.
\end{definition}

\begin{definition}[Étoilé] \label{def: Étoilé}
	ON appelle un domaine $U \subset \Rn$ \textbf{étoilé}, pourvu qu'il existe un $p \in U$ tel que pour tout $x \in U$, le segment de ligne droite $[p,x]$ est inclus dans $U$. De plus on dira que $U$ est \textbf{convexe} pourvu que $U$ est étoilé par rapport à chacun de ses points.
\end{definition}

\begin{lemme}
	Si $U$ est étoilé alors $U$ est simplement connexe.
\end{lemme}

\prf{Soit $\gamma : [0,1] \to U$ un lacet. Soit $p$ dans la définition, alors on pose
	\[
	\Gamma(s,t) = sp + (1-s) \gamma(t), \quad s \in [0,1], \, t \in [0,1]
	\]
	Si $s = 1, \quad \Gamma(1,t) = p$ donc c'est un lacet constant. Et $\Gamma(s,t) \in U$ pour tout $(s,t) \in [0,1]^2$. 
}


\subsection{Green-Riemann et potentiels scalaires}
\begin{proposition}[Partition de l'unité] \label{prop: Partition de l'unité}
	Soit $K \subset \Rn$ un sous-ensemble compact, soit $U_i, V_i$ pour $i = 1, \ldots, N$ des domaines tels que $\overline U_i \subset V_i$, chaque $U_i$ est borné, et 
	\[
	K \subset \bigcup_{i=1}^N U_i
	\]
	Alors, il existe des fonctions $\phi_j \in C^\infty_0(V_j), \quad \phi_j \geqslant 0$ et telles que
	\[
	\forall x \in K, \quad \sum_{j=1}^N \phi_j(x) = 1
	\]
\end{proposition}

\lecture{14 septembre}

\prf{On construit les fonctions $g_i \in C^0_0(V_i)$ avec ${g_i}_{\mid \overline U_i} = 1$ pour tout $i \in \llbracket 1, N \rrbracket$. Pour arriver à $g_i$, on va considérer la composition de la fonction $\dist(x, U_i)$ avec une autre fonction. Ainsi soit $\varepsilon_i = \dist(U_i, V_i^c) = \dist(\overline U_i, V_i^c) > 0$ car $\overline U_i$ est compact, $V_i^c$ fermé, et $\overline U_i \cap V_i^c = \emptyset$. Donc soit $\varepsilon = \min_{i \in \llbracket 1,N \rrbracket} \varepsilon_i$. On définit la fonction suivante
	\[
	\begin{array}{cccc}
		\zeta : & \R_+ & \to & \R \\
		& t & \mapsto & \begin{cases}
			1 & 0 \leqslant t \leqslant \varepsilon/4 \\
			(\varepsilon/2 - t) / (\varepsilon/4) & \varepsilon/4< t < \varepsilon/2 \\
			0 & t \geqslant \varepsilon/2
		\end{cases}
	\end{array}
	\]
	Ainsi il nous suffit de poser $g_i(x) = \zeta(\dist(x, U_i))$. On a $g_i \in C^0(\R)$.
	
	Il reste donc à régulariser les $g_i$. Soit $\psi \in C^\infty_0(B_1(0))$, de plus on peut supposer que $\psi \geqslant 0$ et $\int_{\Rn} \psi(x) \, \mathrm d x = 1$. Pour chaque $\delta > 0$, on définit $\psi_\delta(x) = \delta^{-n} \psi\!\left(\frac{x}{\delta}\right)$. On choisit en particulier $\delta < \varepsilon/4$. On remplace les fonctions $g_i$ par les intégrales paramétrées,
	\[
	\tilde \varphi_i(x) = \int_{\Rn} g_i(x -y) \psi_\delta (y) \, \mathrm dy
	\]
	Il faut vérifier que $\tilde \varphi_i$ satisfait encore $\left.\tilde \varphi_i\right|_{\overline U_i} = 1$, et $\left.\tilde \varphi_i\right|_{V_i^c} = 0$. En effet, si $x \in \overline U_i$, on a 
	\[
	\dist(x-y, U_i) \leqslant \dist(x, U_i) + \norme y \leqslant \dist(x, U_i) + \delta < \varepsilon/4
	\]
	par le choix de $\delta$. Ainsi on a
	\[
	\forall x \in \overline U_i, \quad \tilde \varphi_i(x) = \int_{B_\delta(0)} g_i(x -y) \psi_\delta(y) \, \mathrm d y= \int_{B_\delta(0)} \psi_\delta (y) \, \mathrm dy = 1
	\]
	De la même manière si $x \in V_i^c$ et $y \in B_\delta(0)$ alors 
	\[
	\dist(x-y, U_i) \geqslant \dist(x, U_i) - \norme y \geqslant \varepsilon_i - \varepsilon/4 \geqslant \varepsilon - \varepsilon/4 = \frac{3 \varepsilon}{4}
	\]
	Ainsi $\zeta(\dist(x-y, U_i)) = 0$ si $x \in \overline U_i, ~ y \in V_i^c$ et donc $g_i(x-y) = 0$. Ainsi 
	\[
	\forall x \in V_i^c, \quad \tilde \varphi_i(x) = \int_{B_\delta(0)} g_i(x-y) \, \psi_\delta(y) \, \mathrm d y = 0
	\]
	Ainsi les propriétés des $g_i$ sont conservées sur les $\tilde \varphi_i$ sauf que désormais $\tilde \varphi_i \in C_0^\infty$.
	
	On construit à présent les $\varphi_i$ à partir de $\tilde \varphi_i$ de sorte que $\sum_{i=1}^N \varphi_i (x) = 1$ sur $K$. On pose pour tout $i \in \llbracket 1, N \rrbracket$,
	\[
	\varphi_i = (1 -\tilde \varphi_1)(1 - \tilde \varphi_2) \cdot \ldots \cdot (1 - \tilde \varphi_i)
	\]
	Par récurrence immédiate on vérifie que 
	\[
	1 - \sum_{i=1}^N \varphi_i(x) = (1 - \tilde \varphi_1)(1 - \tilde \varphi_2) \cdot \ldots \cdot (1 - \varphi_N)
	\]
	Et donc si $x \in K \subset \bigcup_{i=1}^N$, alors il existe $j$ tel que $x \in U_j$, donc $\tilde \varphi_j(x) = 1$ et par conséquent $1 - \sum_{i=1}^N \varphi_i(x) = 0$. Ainsi $\sum_{i=1}^N \varphi_i(x) = 1$ sur $K$.
}


\begin{lemme}[Green-Riemann sur des régions de type I ou II] \label{lem: Green-Riemann sur des régions de type I ou II}
Soit $V$ un région de type $I$, c'est-à-dire qu'il existe $h_{1,2} \in C^1([a,b], \R)$ avec $h_1 < h_2$ tel que
\[
V = \{ (x,y) \in \R^2 \mid x \in [a,b], \quad h_1(x) < y < h_2(x)\}
\]
Alors la conclusion de Green-Riemann s'applique
\[
\int_\gamma \bm f \cdot \mathrm d \bm s = \iint_V \rot \bm f \, \mathrm d x \mathrm dy
\]
pour tout $\bm f \in C^1$, où $\gamma$ est une paramétrisation de $\partial V$ entourant $V$ au sens positif.
\end{lemme}

\prf{On calcule explicitement les deux cotés, ainsi on pour le coté droite
	\[
	\iint_V \partial_x f_2 - \partial y f_1 \, \mathrm d x \, \mathrm dy = \underbrace{\int_a^b \int_{h_1(x)}^{h_2(x)} \partial_x f_2 (x,y) \, \mathrm d x \,\mathrm dy}_A - \underbrace{\int_a^b \int_{h_1(x)}^{h_2(x)} \partial_y f_1 (x,y) \, \mathrm d x \, \mathrm dy}_B
	\]
	Ainsi pour $A$ on trouve
	\begin{align*} 
	A &=  \int_a^b \frac{\partial }{\partial x} \left[ \int_{h_1(x)}^{h_2(x)} f_2(x,y) \, \mathrm dy \right] \mathrm dx - \int_a^b h_2'(x) f_2(x, h_2(x)) \, \mathrm dx + \int_a^b h_1'(x) f_2(x, h_1(x)) \, \mathrm dx \\
	&=\int_{h_1(b)}^{h_2(b)} f_2(b,y) \, \mathrm dy - \int_{h_1(a)}^{h_2(a)} f_2(a,y) \, \mathrm dy - \int_a^b h_2'(x) f_2(x, h_2(x)) \, \mathrm dx + \int_a^b h_1'(x) f_2(x, h_1(x)) \, \mathrm dx 
	\end{align*}
		et pour $B$ on a plus simplement,
	\[
	B =\int_a^b \left(\int_{h_1(x)}^{h_2(x)} \partial_y f_1 (x,y) \, \mathrm d x \right)\mathrm dy = \int_a^b f_1(x, h_2(x)) \, \mathrm dx - \int_a^b f_1(x, h_1(x)) \, \mathrm dx
	\]
	Ensuite par concaténation on pose $\gamma_1 = \{(x,h_1(x)) \mid x \in [a,b]\}$, alors
	\[
	\int_{\gamma_1} \bm f \cdot \mathrm d \bm s = \int_a^b \bm f(x, h_1(x)) \cdot \begin{pmatrix}
		1 \\ h_1'(x) \end{pmatrix} \,\mathrm dx = \int_a^b f_1(x, h_1(x)) \, \mathrm dx + \int_a^b h_1'(x) f_2(x, h_1(x)) \, \mathrm dx 
	\]
	De plus on pose $\gamma_2 = \{ (b,y) \mid y \in [h_1(b), h_2(b)]\}$ ainsi
	\[
	\int_{\gamma_2} \bm f \cdot \mathrm d \bm s = \int_{h_1(b)}^{h_2(b)} \bm f(b,y) \cdot \begin{pmatrix} 0 \\ 1 \end{pmatrix} \, \mathrm dy = \int_{h_1(b)}^{h_2(b)} f_2(b,y) \, \mathrm dy 	
	\]
	En posant $-\gamma_3 = \{ (x, h_2(x)) \mid x \in [a,b] \}$ alors
	\[
	\int_{\gamma_3} \bm f \cdot \mathrm d \bm s = - \int_{- \gamma_3} \bm f \cdot \mathrm d \bm s = - \int_a^b f_1(x, h_2(x)) + h_2'(x) f_2(x, h_2(x)) \, \mathrm dx 
	\]
	Et finalement avec $\gamma_4 = \{ (a,y) \mid y \in [h_2(a), h_1(x)]\}$ alors
	\[
	\int_{\gamma_4} \bm f \cdot \mathrm d \bm s = - \int_{h_1(a)}^{h_2(a)} f_2(a, y) \, \mathrm d \bm s
	\]
	Et donc par concaténation on a bien 
	\[
	\int_\gamma \bm f \cdot \mathrm d \bm s = \int_{\gamma_1} \bm f \cdot \mathrm d \bm s + \int_{\gamma_2} \bm f \cdot \mathrm d \bm s + \int_{\gamma_3} \bm f \cdot \mathrm d \bm s + \int_{\gamma_4} \bm f \cdot \mathrm d \bm s = A - B = \iint_V \rot \bm f \, \mathrm d x \mathrm dy
	\]
	D'où le lemme. La preuve les des domaines de types II est analogue.
}
\begin{theorem}[Green-Riemann] \label{thm: Green-Riemann}
	Soit $U \subset \Rn$ un domaine, et soit $V \subset U$ un sous-domaine borné, avec $\overline V \subset U$, borné pour un lacet $\gamma$ régulier, simple, orienté positivement. Alors on a que
	\[
	\int_\gamma \bm f \cdot \mathrm d \bm s = \iint_V \rot \bm f \mathrm d x \, \mathrm dy
	\]
\end{theorem}

\begin{remarque}
	Un lacet $\gamma$ est dit \textbf{simple} $\gamma : [0,1] \to U$ est injective à la condition près $\gamma(0) = \gamma(1)$. Un lacet $\gamma$ est dit \textbf{régulier} si $\gamma \in C^1([0,1], U)$ avec $\gamma'(t) \neq 0$ pour tout $t \in [0,1]$ et $\gamma'(1) = \gamma'(0)$.
\end{remarque}

\prf{Soit $\gamma$ et $V$ comme donnée dans l'énoncé, et $\bm f \in C^1$, alors par régularité de $\gamma$, on a que 
	\[
	\gamma'(t) = \begin{pmatrix} \gamma_1'(t) \\ \gamma_2'(t) \end{pmatrix} \neq 0, \quad \forall t \in I
	\]
	Soit $t_\star \in I$ tel que $\gamma_1'(t_\star) \neq 0$. Ainsi par le théorème des fonctions implicite, il existe un intervalle ouvert $J$ tel que $\gamma_1(t_\star) \in J$, ainsi qu'un autre intervalle ouvert $\tilde J$ tel que $t_\star \in \tilde J$, tel que
	\[
	\forall x \in J, \, \exists ! t \in \tilde J, \quad \gamma_1(t) = x
	\]
	En plus, si on pose $t = t(x) \in C^1(J, \tilde J)$, on peut écrire $y = \gamma_2(t) = \gamma_2(t(x)) := h(x) \in C^1(J)$. Ainsi localement on a 
	\[
	\{ \gamma(x) \mid t \in \tilde J \} = \{ (x, h(x)) \mid x \in J\}
	\]
	
	Pour chaque $t \in I$ on considère le point $\gamma(t)$, et on prend un rectangle $V$ ouvert aligné avec les axes de coordonnées, tel que $V \cap \gamma $ est le graphe d'une fonction par rapport à $x$ ou à $y$. On peut aussi choisir un rectangle ouvert $U$ tel que $\overline U \subset V$ et $\gamma(t) \in U$. On fait ça pour tout $t \in I$, ainsi les $\{U\}$ forment un recouvrement de $\gamma(I)$ par des ouverts. Ainsi par compacité de $\gamma(I)$, d'après Heine-Borel, on peut choisir un nombre fini $\{U_i, V_i\}_{i=i}^N$ tels que $\overline U_i \subset V_i$ et 
	\[
	\gamma(I) \subset \bigcup_{i=1}^N U_i 
	\]
	Pour obtenir un recouvrement de tout $\overline V$. On rajoute encore une paire $U_{N+1}, V_{N+1}$ construit comme suit ; par construit les $U_i$ déjà choisies forment un recouvrement ouvert de $\gamma(I)$, on peut se convaincre qu'il existe $\delta > 0$ tel que
	\[
	\{ x \in V \mid \dist(x, \gamma(I)) < \delta \} \subset \bigcup_{i=1}^N U_i
	\]
	Ensuite on pose,
	\[
	U_{N+1} = \{x \in V \mid \dist(x, \gamma(I)) > \frac\delta2\} \et V_{N+1} = \{x \in V \mid \dist(x, \gamma(I)) > \frac\delta4\}
	\]
	Maintenant qu'on a un recouvrement, on utilise la \autoref{prop: Partition de l'unité} précédente avec $K = \overline V$, et $(U_i, V_i)$, $i \in \llbracket 1, N+1 \rrbracket$. On obtient alors des fonctions
	\[
	\varphi_j \in C^\infty_0 (V_j), \quad \varphi_j \geqslant 0, \quad \sum_{j=1}^{N+1} \varphi_j(x) = 1,\quad \forall x \in \overline V
	\]
	On a alors, pour le champ $\bm f$,
	\[
	\bm f = \sum_{j=1}^{N+1} \varphi_j \bm f
	\]
	Il s'ensuit alors que 
	\[
	\int_\gamma \bm f \cdot \mathrm d \bm s = \int_\gamma \sum_{j=1}^{N+1} \varphi_j \bm f \cdot \mathrm d \bm s = \sum_{j=1}^{N+1} \int_\gamma \varphi_j \bm f \cdot \mathrm d \bm s
	\]
	et de même
	\[
	\iint_V \rot \bm f \, \mathrm dx \, \mathrm dy = \iint_V \rot \sum_{j=1}^{N+1} \varphi_j \bm f \, \mathrm dx \, \mathrm dy = \sum_{j=1}^{N+1} \iint_V \rot(\varphi_j \bm f) \, \mathrm dx \, \mathrm dy
	\]
	Donc il suffit de montrer l'égalité terme par terme. En appliquant le \autoref{lem: Green-Riemann sur des régions de type I ou II} alors 
	\[
	\int_{\partial (V \cap V_j)} \varphi_j \bm f \cdot \mathrm d \bm s = \iint_{V \cap V_j} \rot(\varphi_j \bm f) \, \mathrm d x \, \mathrm dy 
	\]
	Mais comme $\varphi_j \bm f$ est nul en dehors de $V_j$, alors 
	\[
	\iint_V \rot(\varphi_j \bm f) \, \mathrm dx \, \mathrm dy = \iint_{V \cap V_j} \rot(\varphi_j \bm f) \, \mathrm dx \, \mathrm dy
	\]
	De plus le bord de $V \cap V_j$ est constitué de deux morceaux, un morceau de $\gamma$ et un morceau de bord de $V_j$. Or comme $\varphi_j$ est à support compact dans $V_j$, alors elle nulle sur $\partial V_j$ et donc
	\[
	\int_\gamma \varphi_j \bm f \cdot \mathrm d \bm s = \int_{\partial(V \cap V_j)}  \varphi_j \bm f \cdot \mathrm d \bm s
		\]
		On a donc montré que $\int_\gamma \varphi_j \bm f \cdot \mathrm d \bm s = \iint_V \rot(\varphi_j \bm f) \, \mathrm dx \, \mathrm dy$ pour tout $j = 1, 2, \ldots, N$. De plus si $j = N+1$, alors
		\[
		\int_{\gamma} \varphi_{N+1} \bm f \cdot \mathrm ds = 0
		\]
		parce que $\varphi_{N+1}$ n'intersecte pas le lacet $\gamma$. De plus on a
		\[
		\iint_V \rot(\varphi_{N+1} \bm f) \, \mathrm dx \, \mathrm dy = \iint_{\R^2} \rot(\varphi_{N+1} \bm f) \, \mathrm dx \, \mathrm dy = 0
		\]
		car le support de $\varphi_{N+1}$ est un sous-ensemble de $V$. Et donc on a bien l'égalité souhaité
		\[
		\forall j \in \llbracket 1, N+1 \rrbracket, \quad \int_\gamma \varphi_j \bm f \cdot \mathrm d \bm s = \iint_V \rot(\varphi_j \bm f) \, \mathrm dx \, \mathrm dy
		\]
		Ce qui conclut la preuve.	
}

\begin{proposition}
	Soit $U \subset \Rn$ un domaine et $\bm f \in C^1(U, \Rn)$. Supposons que $\rot \bm f = 0$.
	\begin{enumerate}
		\item Soient $\gamma_1, \gamma_2 : [0,1] \to U$ deux chemines en $U$ avec les mêmes extrémités qui sont homotopes, alors
		\[
		\int_{\gamma_1} \bm f \cdot \mathrm d \bm s = \int_{\gamma_2} \bm f \cdot \mathrm d \bm s
		\]
		\item Soit $\gamma : [0,1] \to U$ un lacet contractible, alors
		\[
		\int_\gamma \bm f \cdot \mathrm d \bm s = 0
		\]
		\item Si $U$ est simplement connexe, alors $\bm f$ dérive d'un potentiel, c'est-à-dire
		\[
		\bm f = \nabla g
		\]
	\end{enumerate}
\end{proposition}

\lecture{17 septembre}

